\section{System Model and Security Goals}
\label{sec:model}

This section fixes who participates, what may fail, what is trusted, and which
properties the protocol is meant to give in each FSM state.

\subsection{Validators, Network, and Assumptions}
\label{subsec:threat-model}

\textbf{Validators and network.} We assume partial synchrony: message delay
is unbounded before an unknown global stabilization time (GST) and bounded
after it~\cite{dwork1988consensus}. Let $V$ be the validator set, $w_i > 0$
the voting power of $i \in V$, $W = \sum_{i \in V} w_i$ the total voting
power, and $\operatorname{wt}(S) = \sum_{i \in S} w_i$ the weight of
$S \subseteq V$ (Table~\ref{tab:notation}). The adversary controls a set
$B \subseteq V$ of Byzantine validators, which may deviate arbitrarily and
collude; the other validators $H = V \setminus B$ are honest. As in
CometBFT, a \emph{quorum} is a set whose weight exceeds two thirds of the
total:
\begin{equation}
\operatorname{wt}(B) < \frac{W}{3}, \qquad
\operatorname{wt}(S) > \frac{2W}{3} \ \text{ for every quorum } S.
\label{eq:quorum}
\end{equation}
The seven-validator testbed gives every validator the same power, so
Eq.~\eqref{eq:quorum} allows $f{=}2$ Byzantine validators and a quorum is any
five validators. The adversary cannot forge signatures, find hash collisions,
or subvert the UltraHonk backend's KZG setup. After GST, the \emph{core
consensus network} connects a quorum, as CometBFT itself requires.

\begin{table}[h]
\caption{Notation used in Sections~\ref{sec:model}--\ref{sec:formal}.}
\label{tab:notation}
\centering
\footnotesize
\begin{tabular}{@{}ll@{}}
\toprule
\textbf{Symbol} & \textbf{Meaning} \\
\midrule
$V$, $w_i$, $W$ & validator set, voting power of validator $i$, total voting power \\
$\operatorname{wt}(S)$ & total voting power of the validators in $S$ \\
$B$, $H$ & Byzantine and honest validators \\
$h$, $r$ & consensus height and round \\
$m_i$, $o_i$ & local metric vector and signed observation of validator $i$ \\
$C^{h}$, $\mathsf{Agg}$ & health certificate of height $h$, deterministic aggregation \\
$s^{h}$, $\delta$ & FSM state committed at height $h$, transition function \\
$\tau$ & receipt tolerance \\
\bottomrule
\end{tabular}
\end{table}

\textbf{Validator set.} The validator set is the genesis set and stays fixed
throughout a \textsc{Sovereign} or \textsc{Recovering} interval. Certificate
signatures verify against the genesis validator keys, and the compiled
recovery circuit fixes the set size at seven. Validator-set changes inside a
degraded interval would require proving the set transitions and are outside
this paper.

\textbf{Peripheral layers.} Bitcoin, Celestia, and the peripheral P2P
interfaces may be slow, unavailable, partitioned from some validators, or
adversarially degraded (delayed, censored, eclipsed) for arbitrarily long.
They are not assumed to be Byzantine-fault-tolerant from the validator set's
perspective, but Bitcoin reorganizations deeper than a configurable depth
$K_{\text{DEEP\_FINALITY}}$ are excluded, as in every Bitcoin-anchoring
design (Section~\ref{sec:related}). Honest validators' peripheral
observations may differ from one another at any instant; we do \emph{not}
assume identical views. The FSM transition rests on the health certificate
(Section~\ref{subsec:health-certificate}), which gives every validator the
same aggregate under arbitrary honest disagreement. Whether an honest
sensor reading matches the external layer is a separate question
(Section~\ref{subsec:sensors}).

\textbf{Parameters.} FSM thresholds, hysteresis limits, and
$K_{\text{DEEP\_FINALITY}}$ are genesis parameters shared by every validator,
like the block time or quorum size. A validator running different parameters
simply fails to validate honest proposals; this affects only its own
participation, not safety.

\textbf{Layered adversary.} The adversary may degrade the peripheral
connectivity of \emph{every} honest validator at once (delay Bitcoin
headers, cut the DA bridge, eclipse the peripheral P2P interface), but it
cannot partition the core consensus network below a quorum.
Degrading the periphery therefore changes which FSM state the chain is in; it
does not by itself stop block production.

\textbf{Byzantine power of one third or more.} Beyond $W/3$, no BFT protocol
retains agreement. In \textsc{Anchored} and \textsc{Suspicious}, a
checkpoint that carries the CometBFT block hash lets any Bitcoin observer
detect a fork, because two conflicting histories cannot both match a
Bitcoin-timestamped hash; this is accountability, not prevention. In
\textsc{Sovereign} nothing outside the validator set constrains the chain,
which is why withdrawals are locked in that state
(Section~\ref{subsec:security-semantics}).

\subsection{Security Goals by Mode}
\label{subsec:security-semantics}

We state what each FSM state guarantees, using four properties of the
committed chain $B_1, \dots, B_h$. A \emph{prefix} $B_1, \dots, B_c$ is the
chain up to height $c \leq h$.
\begin{itemize}
  \item $P_{\text{agree}}$: no two conflicting blocks are committed at the
  same height. Assumption: $\operatorname{wt}(B) < W/3$.
  \item $P_{\text{settle}}(c)$: the prefix up to $c$ cannot be rewritten
  without a Bitcoin reorganization deeper than $K_{\text{DEEP\_FINALITY}}$.
  Assumptions: such reorganizations are excluded
  (Section~\ref{subsec:threat-model}), and the Bitcoin checkpoint that
  buries height $c$ commits the chain through $c$ (assumption M2 of
  Theorem~\ref{thm:withdrawal}, Section~\ref{subsec:analytical}).
  \item $P_{\text{avail}}(c)$: the block data of every height up to $c$ is
  retrievable from the DA layer. Assumption: the DA layer is honest and the
  data was published.
  \item $P_{\text{exit}}$: a withdrawal commits only against a chain whose
  settled and available prefix is the one the withdrawal depends on, up to
  the receipt tolerance defined below. It is enforced by the circuit
  breaker.
\end{itemize}

Table~\ref{tab:mode-security} lists which properties hold in each state.

\begin{table}[h]
\caption{Properties by FSM state. $P_{\text{agree}}$ holds in every state
under $\operatorname{wt}(B) < W/3$. $^\dagger$ holds only if the checkpoint
commits the chain through its height and the recovery exit gate includes DA
backfill (Section~\ref{subsec:reanchoring}). The prototype implements both;
the proof condition of the exit gate is qualified in
Section~\ref{subsec:limitations}.}
\label{tab:mode-security}
\centering
\footnotesize
\begin{tabular}{@{}lllll@{}}
\toprule
\textbf{State} & $P_{\text{agree}}$ & $P_{\text{settle}}$ &
$P_{\text{avail}}$ & \textbf{Withdrawals} \\
\midrule
\textsc{Anchored} & Yes & Yes$^\dagger$ & Yes$^\dagger$ & Permitted \\
\textsc{Suspicious} & Yes & Yes$^\dagger$, delayed & Yes$^\dagger$ & Permitted \\
\textsc{Sovereign} & Yes & No & No (validator block stores only) & Locked \\
\textsc{Recovering} & Yes & Pending recovery proof & Backfilling & Locked \\
\bottomrule
\end{tabular}
\end{table}

\textbf{Receipt tolerance.} Let $\tau$ be the receipt tolerance of the round
in which a block commits (Section~\ref{subsec:proposal}). In
\textsc{Anchored} and \textsc{Suspicious}, a valid block at height $h$ whose
layers are reported healthy carries a Bitcoin receipt and a DA receipt that
attest settled and available heights $c_s, c_a \geq h - \tau$.

Theorem~\ref{thm:withdrawal} (Section~\ref{subsec:analytical}) states the conditions under which
a committed withdrawal satisfies $P_{\text{exit}}$. Section~\ref{sec:protocol}
defines the mechanisms, and Section~\ref{sec:formal} states and checks the
properties.
